\section{Exact-content expansion and progression resolution}
\label{sec:content-progression}

\subsection{Content inversion inside the low window}

The content in \eqref{eq:affine-native-content} is
\[
 G_\theta(t)=\gcd(
 \sigma_\theta U_\theta(t),P_\theta).
\]
For \(c\ge1\), exact common-divisor inversion gives
\begin{equation}
 \one_{G_\theta(t)=c}
 =
 \sum_{\kappa\ge1}\mu(\kappa)
 \one_{c\kappa\mid\sigma_\theta U_\theta(t)}
 \one_{c\kappa\mid P_\theta}.
\label{eq:content-exact-projector}
\end{equation}
The sum is finite.

\begin{lemma}[Phase-compatible content inversion]
\label{lem:phase-content-inversion}
For every native key \(\theta\), affine parameter \(t\), and
frequency \(r\),
\begin{align}
 &\one_{G_\theta(t)\le C_X}
 e_{q_X}\!\left(
 -r\varepsilon_\theta
 \frac{M_\theta(t)-n_\theta}{G_\theta(t)}\right)
\notag\\
 &\quad=
 \sum_{1\le c\le C_X}\sum_{\kappa\ge1}
 \mu(\kappa)
 \one_{c\kappa\mid P_\theta}
 \one_{c\kappa\mid\sigma_\theta U_\theta(t)}
 e_{cq_X}\!\left(
 -r\varepsilon_\theta
 (M_\theta(t)-n_\theta)\right).
\label{eq:phase-content-inversion}
\end{align}
\end{lemma}

\begin{proof}
Fix \(c\).  If \(c\kappa\) divides both targets, primitive target
support implies \(c\mid M_\theta(t)-n_\theta\), so the lifted
character on the right is well defined.  Summing \(\mu(\kappa)\)
over divisors of \(G_\theta(t)/c\) vanishes unless
\(G_\theta(t)=c\).  On the surviving term,
\[
 e_{cq_X}\!\left(
 -r\varepsilon_\theta(M_\theta(t)-n_\theta)\right)
 =
 e_{q_X}\!\left(
 -r\varepsilon_\theta
 \frac{M_\theta(t)-n_\theta}{c}\right).
\]
Summing \(c\le C_X\) proves the identity.
\end{proof}

\begin{theorem}[Complete content-expanded low window]
\label{thm:content-expanded-window}
With the notation of \cref{thm:decorated-affine-export},
\begin{align}
 \Low_{K,R}(A_X)
 ={}&
 \sum_{\theta\in\Theta_X}
 \sum_{1\le c\le C_X}
 \sum_{\kappa\ge1}
 \sum_{r\in\Win_X(R)}
 c_{\theta,X}\mu(\kappa)
 \mathfrak m_{K,X}(r)
 \one_{c\kappa\mid P_\theta}
 \notag\\
 &\quad\times
 \sum_{t\in I_\theta}
 \mu(D_\theta(t))\mu(U_\theta(t))
 \chi_\theta(t)W_{\theta,X}(t)
 \one_{c\kappa\mid\sigma_\theta U_\theta(t)}
 \notag\\
 &\quad\times
 e_{cq_X}\!\left(
 -r\varepsilon_\theta
 (M_\theta(t)-n_\theta)\right).
\label{eq:content-expanded-window}
\end{align}
Every sum is finite.  In particular, the right side is one coherent
outer sum; the \((c,\kappa,r)\) terms have not been replaced by their
absolute values.
\end{theorem}

\begin{proof}
Expand \(\Ker_{K,R,X}\) in
\eqref{eq:decorated-affine-export}, apply
\cref{lem:phase-content-inversion}, and interchange finite sums.
\end{proof}

\begin{remark}[Why \(c\) and \(\kappa\) are different keys]
The integer \(c\) is the candidate exact content and determines the
phase modulus \(cq_X\).  The integer \(\kappa\) is a signed
inclusion--exclusion variable.  Replacing \(c\) by
\(b=c\kappa\) in the phase would change the normalized determinant
character.
\end{remark}

\subsection{Solving the native divisibility condition}

Fix an active tuple \((\theta,c,\kappa)\), put
\begin{equation}
 b=c\kappa,\qquad
 g_{\theta,b}=(b,\sigma_\theta),\qquad
 B_{\theta,b}=\frac{b}{g_{\theta,b}}.
\label{eq:progression-B}
\end{equation}

\begin{lemma}[One progression or none]
\label{lem:one-progression}
The native divisibility condition satisfies
\begin{equation}
 b\mid\sigma_\theta U_\theta(t)
 \quad\Longleftrightarrow\quad
 B_{\theta,b}\mid U_\theta(t).
\label{eq:progression-divisibility-reduction}
\end{equation}
Moreover:
\begin{enumerate}[label=(\roman*),leftmargin=2.2em]
\item if \((a_\theta,B_{\theta,b})>1\), there is no solution;
\item if \((a_\theta,B_{\theta,b})=1\), there is one residue
\(\tau_{\theta,b}\pmod{B_{\theta,b}}\) such that every solution is
\[
 t=\tau_{\theta,b}+B_{\theta,b}z,\qquad z\in\Z.
 \]
\end{enumerate}
\end{lemma}

\begin{proof}
Write \(b=gB\) and \(\sigma_\theta=g\sigma'\).  By the definition of
\(g\), \((B,\sigma')=1\).  Therefore
\[
 gB\mid g\sigma' U_\theta(t)
 \quad\Longleftrightarrow\quad B\mid U_\theta(t),
\]
which proves \eqref{eq:progression-divisibility-reduction}.

The inherited primitive relations imply
\((a_\theta,u_\theta)=1\): a common divisor would divide
\(\sigma_\theta u_\theta-a_\theta d_\theta=h_0\), contrary to
\((a_\theta,h_0)=1\).  Hence if a prime divides both
\(a_\theta\) and \(B\), it does not divide \(u_\theta\), and the
congruence \(u_\theta+a_\theta t\equiv0\pmod B\) is insoluble.
When \((a_\theta,B)=1\), the slope is invertible modulo \(B\), giving
one residue class.
\end{proof}

On the soluble branch, define
\begin{align}
 I_{\theta,b}^\ast
 &=
 \{z\in\Z:
 \tau_{\theta,b}+B_{\theta,b}z\in I_\theta\},
\label{eq:progression-resolved-interval}\\
 D_{\theta,b}^\ast(z)
 &=
 D_\theta(\tau_{\theta,b}+B_{\theta,b}z),
\label{eq:progression-D}\\
 U_{\theta,b}^\ast(z)
 &=
 U_\theta(\tau_{\theta,b}+B_{\theta,b}z),
\label{eq:progression-U}\\
 M_{\theta,b}^\ast(z)
 &=
 \ell_\theta v_\theta D_{\theta,b}^\ast(z).
\label{eq:progression-M}
\end{align}
The set \(I_{\theta,b}^\ast\) is an interval for each interval
component already contained in \(\theta\), possibly empty.  Put
\begin{equation}
 \Omega_{\theta,b}
 =
 \ell_\theta v_\theta\sigma_\theta B_{\theta,b}.
\label{eq:progression-phase-slope}
\end{equation}

\begin{proposition}[Resolved phase formula]
\label{prop:resolved-phase}
On a soluble branch,
\begin{align}
 &e_{cq_X}\!\left(
 -r\varepsilon_\theta
 (M_{\theta,b}^\ast(z)-n_\theta)\right)
\notag\\
 &\quad=
 \zeta_{\theta,b,c,r}\,
 e_{cq_X}(-r\varepsilon_\theta\Omega_{\theta,b}z),
\label{eq:progression-resolved-phase}
\end{align}
where
\[
 \zeta_{\theta,b,c,r}
 =
 e_{cq_X}\!\left(
 -r\varepsilon_\theta
 (M_\theta(\tau_{\theta,b})-n_\theta)\right)
\]
has modulus one.
\end{proposition}

\begin{proof}
The moving row satisfies
\[
 M_{\theta,b}^\ast(z)
 =
 M_\theta(\tau_{\theta,b})
 +
 \ell_\theta v_\theta\sigma_\theta
 B_{\theta,b}z.
\]
Insert this identity in the character and separate the constant
phase.
\end{proof}

\begin{corollary}[Fully resolved window]
\label{cor:fully-resolved-window}
The right side of \eqref{eq:content-expanded-window} may be restricted
to tuples satisfying
\[
 b=c\kappa\mid P_\theta,\qquad
 (a_\theta,B_{\theta,b})=1,
\]
and the inner \(t\)-sum replaced exactly by
\begin{align}
 \zeta_{\theta,b,c,r}
 \sum_{z\in I_{\theta,b}^\ast}
 &\mu(D_{\theta,b}^\ast(z))
 \mu(U_{\theta,b}^\ast(z))
 \chi_\theta(\tau_{\theta,b}+B_{\theta,b}z)
\notag\\
 &\times
 W_{\theta,X}(\tau_{\theta,b}+B_{\theta,b}z)
 e_{cq_X}(-r\varepsilon_\theta\Omega_{\theta,b}z).
\label{eq:fully-resolved-inner-sum}
\end{align}
No content or frequency term has been omitted.
\end{corollary}

\begin{proof}
Use \cref{lem:one-progression,prop:resolved-phase} in
\eqref{eq:content-expanded-window}.  Insoluble and empty branches
contribute zero.
\end{proof}
